\section{Keeping the group calculation attached to its books}
\label{sec:group-record-bindings}

The disposal example gives the parent a cash receipt while removing cash
from the consolidated balance sheet. A software implementation must know
whose cash each entry describes. If it posts the removal of the subsidiary's
156 to the parent's legal book, it records a payment the parent never made.
If it leaves the subsidiary's cash in the group after control ends, it
overstates the group's assets. The consolidation calculation therefore needs
its own book and exact references to the legal books it reads.

Call the legal entities Parent-P and Subsidiary-S, and the reporting entity
Group-P. Parent-P's functional currency is USD; Subsidiary-S's is EUR.
Their identifiers remain distinct even if two imported records have the same
display name. Each legal-book snapshot identifies its entity, book, currency,
relative date and signed balances. At acquisition the parent's investment
115.2 and cash payment stay in its own USD book. The subsidiary's opening
cash 20, inventory 140 and payable 60 stay in its EUR book. Group-P's
consolidation entries refer to both snapshots without altering either.
This implements the separation between combining balances and eliminating
the parent's investment against subsidiary equity
\citep[IFRS~10.B86--88]{eu20231803}.

The selected perimeter has equally precise dependencies. Its control
assessment records the investor, investee, substantive rights, variable
returns and ability to use power, together with the interval for which
those findings hold. Ordinary ownership of 80\% and proportionate NCI
measurement are separate inputs. The currency record distinguishes the
subsidiary's functional currency from the group's presentation currency.
A rate of 1.30 is usable here only when its units are USD per EUR and its
date is the closing date R. An inverse quote or a later rate is a different
input, even when its label looks familiar. The relative A, O, R and D
calendar defines this synthetic history; it supplies no evidence of when an
actual investor could have known its facts.

The group needs more state than the initial corporate example. Goodwill,
the parent's translation reserve and NCI have distinct balances, while
profit, OCI and cash movements retain their period histories. The proposed
\texttt{GroupForeignOperation/1} profile supplies those fields under its own
schema lock. It preserves the common request and result fields rather than
silently adding accounts to the closed corporate profile. Loading this
extension would still require an implementation to declare support for it;
there is no automatic conversion of a corporate state into a group state.

Follow the existing history through that representation. At A the group
records the acquisition. At O it translates the opening balances before the
subsidiary's operating entries. At R it finishes cash and goodwill
translation and allocates operating profit. A and O are intermediate
worksheet states; they are not presented as fully closed interim reports.
The R state retains cash 340.8, goodwill 20.8, parent translation reserve
10.4 and NCI 31.2, together with profit 25, OCI 12.6 and the separate cash
exchange effect 7. Closing the income accounts does not erase those movements.

Full disposal and disposal with a retained 20\% interest then start from
separate copies of that same immutable R state. Their branch identities
prevent the second alternative from being posted after the first. Each
disposal proposal identifies its opening state, control change, source books,
cash consideration, retained fair value, rates and accounting method.
Its permission to post is conditional on this synthetic profile and exact
proposal. It is not permission to post to an actual company's books.

In the retained-interest alternative, the parent's receipt is 126.75.
The group removes subsidiary cash 156, producing investing cash
$126.75-156=-29.25$. The retained investment 42.25 is a separate noncash
asset. Both disposal alternatives remove the entire NCI balance 31.2,
including its translation reserve 2.2. Only the parent's reserve 10.4 is
recycled into the disposal gain. These are checks on the meaning of the
journal, beyond the fact that its debits equal its credits
\citep[IFRS~10.B98--99; IAS~21.48A--B; IAS~7.39--42]{eu20231803}.

The report reads one branch and one dependency view. In both alternatives,
period profit is 58.8, of which 53.8 belongs to the parent and 5 to NCI.
Period OCI is 2.2, attributable to NCI. Ending NCI is nevertheless zero,
because the disposal removes its accumulated balance. The equity movement
must therefore retain both its introduction and removal:
\[
 300+24+58.8+2.2-31.2=353.8.
\]
Reading only the closing NCI account would lose the subsidiary's contribution
while it was controlled. Reading only final equity would also miss the
difference between full-sale cash 353.8 and retained-interest cash 311.55
plus investment 42.25.

Publication follows the common accounting transaction rule. The request
identifies the group owner and expected opening revision; its business-event
identity includes the branch. Before publishing, an implementation must
validate the source and currency bindings, posting authorization, ordered
journals and independently specified state changes. It then publishes the
closing state, journal, economic delta, receipt and new group head together.
The parent and subsidiary legal-book owners remain untouched.

An acknowledgement can be lost after that publication. Retrying the exact
original request then finds the existing receipt and has no second effect,
even though its original opening revision is now old. If publication never
occurred, the same request can commit once. Changing the request identifier,
budget or expected head changes the fingerprint and cannot masquerade as
that retry. A newly proposed event with an outdated opening is rejected.
The documentary packet preserves both complete acknowledgement histories
and an invalid partial-publication example; later runtime tests must exercise
those failures in the chosen storage system.

A corrected exchange-rate source creates a new evidence view. It makes the
old group's calculation unsuitable for a request using the corrected view,
but does not overwrite the old books or remove the old explanation.
Unknown control or an unknown rate instead returns the missing-input
obligation, with no fabricated zero amount. The companion group-binding
specification supplies exact records and exchanges for these cases.
Actual control, valuation, tax, policy alignment and subsequent accounting
for the retained investment retain their separate evidence requirements.
